\uuid{Qlag}
\exo7id{1727}
\titre{exo7 1727}
\auteur{legall}
\organisation{exo7}
\datecreate{1998-09-01}
\isIndication{false}
\isCorrection{false}
\chapitre{Réduction d'endomorphisme, polynôme annulateur}
\sousChapitre{Applications}
\cpge{mpsi-mp2i.calcul-matriciel-systemes.e, pcsi.determinants.c, ptsi.determinants.a, mpsi-mp2i.matrices-applications-lineaires.a, pcsi.matrices-applications-lineaires.a, ptsi.matrices-applications-lineaires.a, mp-mpi.reduction-endomorphismes.b, pc.reduction.a, psi.reduction-endomorphismes.a, mp-mpi.reduction-endomorphismes.d, pc.reduction.c, psi.reduction-endomorphismes.c, pt.reduction-endomorphismes.c, mp-mpi.reduction-endomorphismes.g, bcpst2.valeurs-propres-vecteurs-propres.b}
\module{Algèbre}
\niveau{L2}
\difficulte{}

\contenu{
\texte{
Soit $  A\in M_n(\mathbb{K} )  .$ On note $  C(A)=\{ B \in M_n(\mathbb{K} )  ;
AB=BA\}   .$
}
\begin{enumerate}
    \item \question{On suppose que $  A  $ a des valeurs propres simples. Montrer que les propri\'et\'es suivantes sont
\'equivalentes~:
\vskip1mm
\hskip1mm i) $  B\in C(A)  .$
\vskip1mm
\hskip1mm ii) $  B  $ a une base de vecteurs propres en commun avec $  A  .$
\vskip1mm
\hskip1mm iii) Il existe $  P\in \mathbb{K} _{n-1}[X]  $ tel que $  B=P(A)  .$
\vskip1mm
\hskip1mm iv) Il existe $  P\in \mathbb{K} [X]  $ tel que $  B=P(A)  .$
\vskip1mm}
    \item \question{On suppose que $  n =3  $ (pour simplifier) et que $  A  $ est diagonalisable
avec une valeur propre double. D\'eterminer $  C(A)  .$}
\end{enumerate}
}
